% Exact evaluation of the already proved finite QGT error at original m=1.
% Provider: QUANTITATIVE_ORIGINAL_GAMMA_RETURN.tex, QGT1--28.
\section{The finite original \(m=1\) return bound at the \(q\) scale}
\label{sec:QGQ}

Retain the same original measure, determinants, source products, and
constants $A,M,\mathcal C_\Gamma$ of QGT1--28. At the original value
$m=1$, the packet is exactly
\[
 k=4l+1,\qquad q=(4l+2)^2,\qquad n=q/2,\qquad l\ge1.
 \tag{QGQ1}
\]
The explicit two finite nonnegative bounds in QGT16 and QGT25 have
the following limits on this packet:
\[
 \frac{\mathcal U_{q,l}}q\longrightarrow\frac A{\sqrt2},
 \qquad
 \frac{\mathcal D_{q,l}}q\longrightarrow
                  \frac M8+\frac{\log2}{4}.
 \tag{QGQ2}
\]
For the first limit use the exact identity
\[
 \frac{\mathcal U_{q,l}}q
 = A\frac l{\sqrt q}
          \left(\frac3{\sqrt2}+\sqrt{\frac12+\frac1q}\right).
\]
Here $l/\sqrt q=l/(4l+2)\to1/4$, so the limit is
$A(1/4)(4/\sqrt2)$, as displayed.

For the second limit, extract from the actual QGT25 sum the two
terms $2Ml^2$ and $4l^2\log2$. Their quotients by $q$ tend to
$M/8$ and $(\log2)/4$, since $l^2/q\to1/16$.
Every other summand tends to zero after division by $q$.
Here is the term-by-term justification using the original summands.
The terms $4Ml$, $l|\mathcal L_0|$, and
$2l|\log(4/\pi)|$ are fixed constants times $l$; their quotients
by $q$ tend to zero. The actual rank correction satisfies
\[
 0\le2(n+1)l\log(1+1/n)
       \le2(n+1)l/n\le4l,
\]
because $\log(1+u)=\int_0^u(1+t)^{-1}dt\le u$ and $n\ge1$.
The two next rational terms of QGT25 have sizes $O(l)$ and $O(1)$,
respectively:
\[
 \frac{8l^3+l}{6q},\qquad\frac{l^4+l^2/4}{q^2}.
\]
To verify these orders directly, substitute
$q=l^2(4+2/l)^2$; the remaining denominators are bounded below
by positive constants for $l\ge1$.
The two nonnegative product bounds have orders
$p^-_{q,l}=O(l^{-1})$ and $p^+_{q,l}=O(l^{-1})$.
For $p^-$ its two terms have orders $l^{-1},l^{-3}$.
For $p^+$ its three displayed terms have orders
$l^{-2},l^{-4},l^{-1}$, respectively. Each follows by the same
literal substitution in QGT21; the maximum therefore also has order
$l^{-1}$.
The remaining three displayed polynomial low terms
\[
 \frac{l^2}{q},\qquad
 \frac{l(16l^2-1)}{48q^2},\qquad
 \frac{l^2(8l^2-1)}{48q^3}
\]
have orders $1,l^{-1},l^{-2}$, respectively.
Finally the actual LER bounds retained in QGT25 give
$0<K_{q,l}\le3l/(32q^3)=O(l^{-5})$,
$T_{q,l}=O(l^{-3})$, and $E_{q,l}=O(l^{-9})$.
These are all the terms of QGT25, and each of these remaining terms
divided by $q\ge16l^2$ tends to zero. This proves the second
limit in QGQ2.

Apply the finite original inequality QGT26 for every integer $l$,
divide it by the same $q$, and use both proved limits in QGQ2.
The result is the quantitative statement
\[
 \boxed{\displaystyle
 \limsup_{l\to\infty}
 \left|\frac{W_k-lq\mathcal C_\Gamma}{q}\right|
 \le\frac A{\sqrt2}+\frac M8+\frac{\log2}{4}
 \qquad(m=1).}
 \tag{QGQ3}
\]
All constants in this inequality are the finite constants of the
proved original-source QGT calculation. The two limits in QGQ2
evaluate its explicit error bounds. Inequality QGQ3 alone does not
determine a limit for the actual expression inside its absolute value;
the exact centre $\mathcal C_{q,l}$ in QGT15 and the signed finite
interval QGT22 remain available for that continuing calculation.
